% !TeX program = lualatex
% Copyright (C) 2026 IEEE UFG Computer Society contributors.
% Distributed under the LaTeX Project Public License 1.3c or later.
% See LICENSE and development/MANIFEST.md.
%
% MODELO DE ESTATUTO
% Substitua os textos de orientação por disposições aprovadas pela organização.
% A numeração de capítulos e artigos é automática e contínua.
\documentclass[
  paper-size=a4,
  font-size=11pt,
  two-sided=true,
  title-page=true,
  inner-margin=22mm,
  outer-margin=22mm,
  top-margin=24mm,
  bottom-margin=28mm,
  color-model=rgb
]{ieeeufgcs}
\usepackage[brazil]{babel}
\usepackage{ieeeufgcs-estatuto}

\IEEECSSetup{
  document-type={Estatuto},
  title={Estatuto do Capítulo Estudantil},
  running-title={Estatuto},
  institution={Capítulo Estudantil da IEEE Computer Society na UFG},
  author={Órgão responsável pelo documento},
  date={\today},
  location={Goiânia},
  identifier={Minuta para revisão}
}

\begin{document}
\maketitle

\EstatutoCapitulo{Disposições gerais}

\EstatutoArtigo
Defina a denominação, a natureza e a finalidade da organização.

\EstatutoArtigo
Indique as normas e os órgãos aos quais a organização está vinculada.

\EstatutoCapitulo{Membros}

\EstatutoArtigo
Defina as condições de ingresso e as categorias de membros.

\begin{EstatutoIncisos}
  \item Descreva a primeira categoria e seus requisitos.
  \item Descreva a segunda categoria e seus requisitos.
\end{EstatutoIncisos}

\EstatutoArtigo
Estabeleça os direitos e os deveres dos membros.

\EstatutoParagrafoUnico
Acrescente uma regra complementar, se necessária.

\EstatutoCapitulo{Organização e funcionamento}

\EstatutoArtigo
Descreva os órgãos, os cargos, os mandatos e suas competências.

\EstatutoParagrafo
Especifique uma condição relativa ao artigo.

\EstatutoCapitulo{Disposições finais}

\EstatutoArtigo
Indique como o estatuto pode ser alterado e como tratar casos omissos.

\end{document}
